\chapter{Design} \label{chap:design}
This chapter will cover the design of a TAPS System. 
This includes the decisions which have affected
CTaps, both for the internal architecture and the external API.
For CTaps the overarching design goal was to conform to the API described in RFC~9622~\cite{rfc9622},
while exploring how to best expose as much QUIC functionality as possible
through the TAPS API.

\section{Choice of Scope}\label{sec:design-scope}
CTaps currently features all major aspects of the TAPS API except for Rendezvous.
Internally it supports  candidate gathering,
candidate racing, TCP, UDP and QUIC.

The following API functions are not implemented:
protocol-specific endpoints, hop limits, multicast endpoints,
resolving Preconnections before initiation,
adding endpoints to an established Connection, 
Connection limits on Listeners,
partial message receives, partial message sends and batching message sends.

All TAPS Transport and Message Properties are present in the CTaps API,
but not all are supported internally.
Unsupported Connection Properties are ignored and have
no effect on the underlying Connection.
Unsupported Selection
Properties still affect candidate gathering, but do not
affect the underlying Connection when selected.
For example, setting \mono{perMsgReliability} to Require will filter
out protocols which have this property set to Prohibit, but the application
will not be able to transfer individual messages as reliable, since this
is not supported in CTaps. 
All Message Properties except for \mono{final} and \mono{safelyReplayable} are ignored.
CTaps only considers Selection Properties of type Preference when
performing candidate gathering.

CTaps allows applications to set all Security Parameters described in RFC~9622,
except for session cache options and pre shared keys.
In addition, CTaps does not support trust callbacks. Only parameters
needed for TLS 1.3 and session resumption have any effect.

For internal TAPS features, CTaps has priortized the most distinct,
namely candidate gathering and racing. This allows a single
client to be able to connect to multiple different endpoints using different
protocols. CTaps does not implement any internal cache or system policy. Additionally,
Listeners in CTaps are only able to listen on a single protocol and a single Local
Endpoint.

CTaps does not present all optional parameters from RFC~9622,
lacking support for timeout when initiating Connections,
minimum incomplete length and maximum length when receiving.

Overall, CTaps has prioritized exposing features in the API even
when not supported internally. The implemented functionality is
nonetheless sufficient to demonstrate the key benefits of a TAPS System

\section{Choice of Dependencies}
\subsection{libuv} \label{sec:dependency-libuv}
To implement asynchrony in CTaps, the event loop model was chosen over multithreading.
This was done for simplicity's sake, as it could avoid shared state.
Splitting control flow across callbacks was not a concern for CTaps, as
these callbacks largely mirrored the event-based API of TAPS.

When selecting a specific event loop library, libuv presented
several advantages. Among these were built-in protocol handles for TCP
and UDP, in addition to being a mature library which is already
used in production systems such as Node.js~\cite{nodejsEventLoop}.
It is also the event loop engine behind
NEAT~\cite{neatCode}, which influenced the decision of CTaps.

Internally, libuv provides the asynchronous interface
which CTaps uses to interact with the network. It additionally powers
the main event loop which applications run by
invoking \mono{ct_start_event_loop()}.
CTaps also uses the timer handle provided by libuv in order to know when to query Picoquic,
according to the sans-I/O model described in \cref{sec:picoquic}.

\subsection{Picoquic}
Since QUIC is not supported natively by libuv, an external library
was needed to add QUIC support to CTaps, with the largest
priority being ease of use.

Picoquic was chosen because
of its simplicity and was
prioritized over larger, more performance-oriented libraries such as
MsQuic~\cite{MsQuic}.

A large advantage of Picoquic is that it features an optional, built-in
event loop. This allowed the experiments in \cref{chap:eval} to be easily performed with a Picoquic
client completely independent of CTaps.

\subsection{GLib} \label{sec:glib-design-reason}
CTaps uses GLib~\cite{GLibGitLab} for implementations of
hash maps, linked lists, trees and arrays.
GLib was added as a dependency since C has a small standard library.
If CTaps were to implement all the data structures internally
it would create a large surface area liable for bugs.

GLib itself was chosen because it is a large and mature library
which is featured in several other well known software libraries,
such as GTK.

\subsection{log.c}
CTaps uses log.c~\cite{logc} to create structured logs, which
can be easily filtered by level.
It was chosen because of its simplicity
and pure C implementation, making it easy to integrate
with CTaps.

\section{Architecture of CTaps} \label{sec:ctaps-arch-design}
Internally in CTaps, the most important architectural concept is the
protocol interface. This is an interface used
to interact with underlying transport-layer protocols, which ensures that CTaps is extensible.
Common logic such as state tracking
is extracted to the higher level concept of the Socket Manager. The Socket Manager is an internal layer 
responsible for multiplexing, maintaining the lifetime of network sockets and communicating
between Connections and the transport protocol.

An overview of the architecture of CTaps can be seen in~\cref{fig:ctaps-arch}.
This displays how the Socket Manager
acts as an intermediary between the CTaps API and the
protocol interface. The unique consideration of QUIC can
also be seen, requiring Picoquic as an external library due to lack
of support in libuv.

\begin{figure}[htbp]
\centering
\begin{tikzpicture}[
    node distance=0.5cm,
    layer/.style={
        draw, 
        line width=0.8pt, 
        minimum width=8cm, 
        minimum height=1.2cm, 
        align=center, 
        fill=gray!10,
        font=\sffamily\bfseries
    },
    sublayer/.style={
        draw, 
        line width=0.6pt, 
        minimum width=3.8cm, 
        minimum height=1cm, 
        fill=white,
        font=\sffamily\small
    }
]

% 1. User Application Layer
\node [layer, fill=blue!10] (app) {Application};

% 2. CTaps API Layer
\node [layer, below=of app] (api) {CTaps API\\ \mdseries\small (TAPS objects, Callback registration)};

% 3. Socket Manager / Abstraction
\node [layer, below=of api] (manager) {Socket Manager\\ \mdseries\small (Owns network socket)};

% 3. Socket Manager / Abstraction
\node [layer, below=of manager] (interface) {Protocol Interface\\ \mdseries\small (Interface to network protocols)};

% % 4. Implementation Engines (Split)
\node [sublayer, below=0.5cm of interface.south west, anchor=north west] (uv) {libuv \mdseries\small (TCP/UDP)};
\node [sublayer, below=0.5cm of interface.south east, anchor=north east] (pico) {Picoquic \mdseries (QUIC)};

\node [sublayer, below=0.5cm of pico.south east, anchor=north east] (libuvraw) {libuv \mdseries (UDP socket)};

% % 5. OS Layer
\node [layer, below=1cm of libuvraw.south west, fill=gray!25] (os) {Network Stack};

% Connectors
\draw[->, thick] (app) -- (api);
\draw[->, thick] (api) -- (manager);
\draw[->, thick] (manager) -- (interface);
\draw[->, thick] (interface.south) -- ++(0,-0.15) -| (uv);
\draw[->, thick] (interface.south) -- ++(0,-0.15) -| (pico);
\draw[->, thick] (uv) -- (uv |- os.north);
\draw[->, thick] (pico) -- (pico |- libuvraw.north);
\draw[->, thick] (libuvraw) -- (libuvraw |- os.north);
\end{tikzpicture}
\caption[CTaps architecture]
{
An overview of the architecture of CTaps, featuring the Socket Manager as a
mediator between the objects presented to applications and the protocol interface.
}
\label{fig:ctaps-arch}
\end{figure}

\subsection{Socket Manager} \label{sec:socket-manager-design}
The Socket Manager is an internal structure in CTaps.
It has three primary responsibilities:
acting as an asynchronous intermediary between the API and the
protocol interface, managing the lifetimes of the underlying sockets
and demultiplexing incoming UDP datagrams.

Since UDP is connectionless, the Socket
Manager maintains an internal hash table to map
incoming UDP datagrams to Connections.
This is done based on the remote socket address.

In order to make each protocol interface stateless,
all protocol data is stored in the associated Socket Manager. Each Socket
Manager is associated with a single protocol interface, but
each protocol interface can be associated with many Socket Managers.
Any state required by the protocol is allocated on the heap and is
stored in void pointers contained in the Socket Manager.
This allows each protocol interface to be global singleton,
eliminating the need to instantiate them per connection.
An example of the TCP singleton utilizing state stored in the Socket Manager
can be seen in \cref{lst:stateless-tcp-example}.

\begin{listing}[htbp]
\begin{minted}{C}
int tcp_close(ct_connection_t* connection) {
    log_info("Closing TCP connection: %s", connection->uuid);

    ct_tcp_socket_state_t* socket_state = connection->socket_manager->internal_socket_manager_state;
    uv_close((uv_handle_t*)socket_state->tcp_handle, on_libuv_close);
    return 0;
}
\end{minted}
\caption[TCP close implementation]{
Example of the TCP interface accessing Socket Manager state when closing a Connection.
}
\label{lst:stateless-tcp-example}
\end{listing}

Since several functions within the protocol interfaces
are asynchronous, the communication between the Socket Manager and the protocol interface also has to be asynchronous.
The Socket Manager features internal
callbacks which are invoked by the protocol
interface at appropriate times, for example to notify
the Socket Manager that a Connection has been closed after teardown
has completed. This callback structure can be seen in~\cref{lst:socket-manager-callbacks}.

For memory-critical paths, Connections exclusively communicate
with the underlying protocol interface through the Socket Manager, in order
to mitigate potential memory leaks.

An example of such a memory-critical internal flow in CTaps can be seen in
\cref{fig:async-close-flow}, showing the asynchronous flow of
graceful connection teardown for a Connection using a shared
Socket Manager. Here, a single Connection in a Connection Group
is closed, and the Socket Manager invokes the application callback
asynchronously.

The figure also shows why it is crucial that the Socket Manager
keeps the underlying socket open, since it is still used by the
remaining Connections.

\definecolor{ConnectionGroupTextColor}{RGB}{40, 61, 41}
\begin{figure}[htbp]
\centering
\begin{tikzpicture}[font=\sffamily,
    block/.style={rectangle, draw, text centered, rounded corners, fill=white, font=\bfseries},
    client/.style={rectangle, draw, text centered, fill=white, font=\small},
    arrow/.style={-Stealth, thick},
    asyncArrow/.style={arrow, dashed, red!70!black},
    label/.style={font=\scriptsize\itshape}
    ]

% Place connections first
\node[client] (Conn1) at (0,0) {Connection 1};
\node[client, below=0.3cm of Conn1] (Conn2) {Connection 2};
\node[client, below=1.2cm of Conn2] (ConnN) {Connection N};

% SM centred vertically on the connections group
\node[block, align=center] (SM) at
    ($(Conn1.north east)!0.5!(ConnN.south east) + (4,0)$)
    {Socket Manager};

\node[block, right=of SM] (Prot) {Protocol Interface};

\draw[arrow] (Conn1.north east) -- node[above, sloped, label, yshift=3pt] {close()} (SM.north west);
\draw[arrow] (Conn2.east) -- (SM.west);
\draw[arrow] (ConnN.east) -- (SM.south west);

\node[font=\Huge, text=gray] at ($(Conn2.south)!0.4!(ConnN.north)$) {$\vdots$};

\begin{scope}[on background layer]
    \draw[rounded corners, fill=blue!5!white, draw=blue!20]
        ($(Conn1.north west)+(-0.3, 0.4)$) rectangle ($(ConnN.south east)+(0.3,-0.3)$);
    \node[font=\normalsize\bfseries, above=0.4cm of Conn1, blue!70!black, align=center]
        {Connections};
    \node[fit=(Conn1)(Conn2), rounded corners, fill=green!5!white,
          draw=green!20, inner sep=0.13cm] (GreenBox1) {};
    \node[fit=(ConnN), rounded corners, fill=green!5!white,
          draw=green!20, inner sep=0.13cm] (GreenBoxN) {};
\end{scope}
\node[font=\tiny\bfseries, text=ConnectionGroupTextColor, above=-1.2mm of GreenBox1] {Connection Group 1};
\node[font=\tiny\bfseries, text=ConnectionGroupTextColor, above=-1.2mm of GreenBoxN] {Connection Group N};

% Request Flow
\draw[arrow]
    ([yshift=-0.5mm]SM.north east)
    -- node[above, sloped, label, align=center] {Close connection}
    ([yshift=-0.5mm]Prot.north west);
\draw[asyncArrow] 
    ([yshift=0.5mm]Prot.south west)
    -- node[below, label, text=red!70!black, align=center] {Socket Manager cb}
    ([yshift=0.5mm]SM.south east);
\draw[asyncArrow]
    ([yshift=4pt]SM.west)
    -- node[below, sloped, label, text=red!70!black, align=center] {User cb}
    ([yshift=3pt]Conn1.east);
\end{tikzpicture}
\caption[Flow of internal callbacks in CTaps]{Diagram of the asynchronous communication
between the CTaps Socket Manager and the protocol interface when an application invokes \mono{close()}.}
\label{fig:async-close-flow}
\end{figure}

\begin{listing}[htbp]
\begin{minted}{C}
typedef struct ct_socket_manager_callbacks_s {
void (*closed_connection)(ct_connection_t* connection);
void (*aborted_connection)(ct_connection_t* connection);
void (*establishment_error)(ct_connection_t* connection);
void (*connection_ready)(ct_connection_t* connection);
void (*message_sent)(ct_connection_t* connection, ct_message_context_t* message_context);
void (*message_send_error)(ct_connection_t* connection,
        ct_message_context_t* message_context, int reason_code);
void (*socket_closed)(struct ct_socket_manager_s* socket_manager);
void (*connection_received)(ct_listener_t* listener, ct_connection_t* connection);
void (*listener_ready)(ct_listener_t* listener);
void (*closed_listener)(struct ct_socket_manager_s* socket_manager);
} ct_socket_manager_callbacks_t;
\end{minted}
\caption[Socket Manager internal callbacks]{The internal callbacks used for communication between
the Socket Manager and the protocol interface.}
\label{lst:socket-manager-callbacks}
\end{listing}

\subsection{Protocol Interface}
In order to make CTaps design as modular as possible,
it uses a generic interface to interact with the underlying transport-layer protocols.
This interface abstracts away the implementation details of the individual
protocols and allows the Socket Manager to not have to implement protocol
specific logic, making it far easier to extend support to more protocols in
the future.

An excerpt of the protocol interface can be seen in
\cref{lst:protocol-impl-struct}. The \mono{supports_alpn} flag
and the Selection Properties are used
to encode the capabilities of the protocol. The encoded capabilities 
allow CTaps to filter and sort protocols according to the
Transport Services desired by the application, while
the implemented function pointers present a generic
interface to protocol-specific behavior. The full generic
protocol interface can be seen in \cref{lst:impl-struct}.

\begin{listing}[htbp]
\begin{minted}{C}
typedef struct ct_protocol_impl_s {
    ///< Protocol name (e.g., "TCP", "UDP", "QUIC")
    const char* name;
    ///< Protocol enumeration value
    ct_protocol_enum_t protocol_enum;
    ///< True if protocol supports ALPN negotiation
    bool supports_alpn;
    ///< Properties supported by this protocol
    ct_selection_properties_t selection_properties;
    /** @brief Initialize a new connection using this protocol. */
    int (*init)(ct_connection_t* connection);
    /** @brief Send a message over the protocol. */
    int (*send)(ct_connection_t*, ct_message_t*, ct_message_context_t*);
    /** @brief Start listening for incoming connections. */
    int (*listen)(struct ct_socket_manager_s* socket_manager);
    // ... Remaining protocol interface function pointers ...
} ct_protocol_impl_t;
\end{minted}
\caption[Protocol interface excerpt]{Excerpt from the protocol interface.}
\label{lst:protocol-impl-struct}
\end{listing}

An example of how the Socket Manager interacts with the protocol interface
can be seen in \cref{lst:example-sock-man-to-protocol}. Connections
invoke \mono{ct_socket_manager_close_connection()}, which is then passed
through the Socket Manager into the protocol interface.

\begin{listing}[htbp]
\begin{minted}{C}
int ct_socket_manager_close_connection(ct_connection_t* connection) {
    assert(connection);
    ct_socket_manager_t* socket_manager = connection->socket_manager;
    log_debug("Socket manager: Closing attached connection: %s", connection->uuid);
    int rc = socket_manager->protocol_impl->close_connection(connection);
    if (rc) {
        log_error("Error from protocol when closing connection: %s", connection->uuid);
        return rc;
    }
    return 0;
}
\end{minted}
\caption[Socket Manager generic protocol API usage]
{
    Example of the access pattern between the Socket Manager and the protocol
    interface. The Connection is not considered to be closed until the
    \mono{closed_connection()} callback is invoked by the protocol interface.
}
\label{lst:example-sock-man-to-protocol}
\end{listing}


\section{API Design}
The external API of CTaps is object-oriented and asynchronous.
CTaps hides internal state from applications and only exposes getters and setters
to specific fields.

Example client and server applications utilizing the CTaps
API can be seen in \cref{app:ex-ctaps-apps}.

\subsection{Object Orientation}
Although C is not an object-oriented language, object-orientation can still
be emulated through the use of structs and dedicated functions.
This is commonly done to achieve a clear separation of concerns.
An example of this separation of concerns in the CTaps API
can be seen with the Remote Endpoint object in~\cref{lst:remote-endpoint}.
This object contains all the context related to Remote Endpoints, which then is used
in Connections and Listeners. The object has dedicated functions to operate
on it, and therefore encapsulates all responsibilities related to Remote Endpoints.

\begin{listing}[htbp]
\begin{minted}{C}
typedef struct ct_remote_endpoint_s {
    ///< Port number set by user
    uint16_t port;
    ///< Service name (e.g., "https") or NULL
    char* service;
    ///< Hostname for DNS resolution or NULL
    char* hostname;
    ///< Resolved socket address
    struct sockaddr_storage resolved_address;
} ct_remote_endpoint_t;

// Example functions
ct_remote_endpoint_t* ct_remote_endpoint_new(void);
int ct_remote_endpoint_with_ipv4(
        ct_remote_endpoint_t* remote_endpoint,
        in_addr_t ipv4_addr
);
void ct_remote_endpoint_with_port(
    ct_remote_endpoint_t* remote_endpoint,
    uint16_t port
);
void ct_remote_endpoint_free(
    ct_remote_endpoint_t* remote_endpoint
);
\end{minted}
\caption[Remote Endpoint structure]
{The internal definition of a Remote Endpoint in CTaps, with example functions from the CTaps API.}
\label{lst:remote-endpoint}
\end{listing}

CTaps presents objects as consistently as possible
both internally and externally.
An example of this is that in CTaps every Connection belongs to a Connection Group,
even if it has not been used for cloning.
This means that Connection Group-oriented operations are always valid
in CTaps. It also allows CTaps to avoid internal checks for whether a Connection was
part of a Connection Group or not.

In order to hide internal state from applications, the public CTaps API only
features forward declarations of objects. This means that the application
cannot interact directly with the struct members and instead must
interact with the objects through a list of functions exposed by CTaps,
which are used to allocate a struct on the heap and to modify the object using
its pointer, a pattern which can be seen in~\cref{lst:using-remote-endpoint}.

\begin{listing}[htbp]
\begin{minted}{C}
    // Step one allocate object on the heap
    ct_remote_endpoint_t* remote_endpoint = ct_remote_endpoint_new();
    // Check for allocation error
    if (!remote_endpoint) {
      return -1;
    }

    // Then modify it
    ct_remote_endpoint_with_ipv4(
        remote_endpoint,
        inet_addr("127.0.0.1")
    );
    ct_remote_endpoint_with_port(remote_endpoint, 1234); // example port

    // ... Then use it in a Preconnection for example...


    // And finally free it when done
    ct_remote_endpoint_free(remote_endpoint);
\end{minted} 
\caption[Remote Endpoint lifetime snippet]
{
    A snippet displaying how an application could manage
    the lifetime of a Remote Endpoint. Allocating on the
    heap and freeing when no longer needed.
}
\label{lst:using-remote-endpoint}
\end{listing}

This pattern was chosen to make it clear exactly which fields are
controllable by the application and which fields are only used internally in CTaps.
This pattern also allows CTaps to differentiate between read-only Connection Properties and writable
Connection Properties as described in the TAPS API~\cite[sec. 8.1.11]{rfc9622}.

\subsection{Asynchrony} \label{asynchrony}
Being asynchronous and built around an event loop has
several implications for the external CTaps API, with the most important being
that all CTaps functions return immediately and never block. If the invocation
initiates an asynchronous piece of work, then results are returned
through a callback.

Registering callbacks in CTaps is done as close
to the relevant event as possible. For example, Connection-related callbacks 
are registered when the Preconnection is initiated. An example of flow
and the registration of callbacks can be seen in~\cref{lst:preconnection-init-flow}.
Not all callbacks have to be registered, although not registering important
callbacks such as \mono{ready()} or \mono{closed()} can lead 
to applications missing important events.

\begin{listing}[htbp]
    \begin{minted}{C}
ct_preconnection_t* preconnection = ct_preconnection_new(
    NULL,
    0,
    &remote_endpoint,
    1,
    transport_properties,
    NULL);

ct_connection_callbacks_t connection_callbacks = {
    .establishment_error = on_establishment_error,
    .ready = send_message_and_receive,
    .per_connection_context = &test_context,
};

int rc = ct_preconnection_initiate(preconnection, &connection_callbacks);
    \end{minted}
    \caption[Preconnection initiation flow]{Example flow registering callbacks when initiating a Preconnection.}
    \label{lst:preconnection-init-flow}
\end{listing}

Since the results of asynchronous operations in CTaps are available through callbacks,
applications need a way to connect these results to state from outside the
callback. This could be achieved through global variables, but CTaps 
also allows applications to pass a void pointer when registering callbacks.

When initiating a Connection, the void pointer will be stored in the Connection
structure, and it can be fetched when the Connection is passed
to a callback. An example of this pattern can
be seen in \cref{lst:conn-ready-context}.

\begin{listing}[htbp]
  \begin{minted}{C}
void on_connection_ready_increment(ct_connection_t* connection) {
    your_custom_struct_t* context = (your_custom_struct_t*)ct_connection_get_callback_context(connection);
    context->num_ready_connections++;
}
  \end{minted}
  \caption[Connection callback context]
{
    Example of accessing user context in a Connection ready callback.
}
  \label{lst:conn-ready-context}
\end{listing}

Connection and Listener
objects are created asynchronously in CTaps as
both of these objects require asynchronous operations
to instantiate.
Creating Connections is asynchronous
because candidate gathering is asynchronous, but
also because creating a Connection may require a handshake to be performed
over the network. Creating Listeners is asynchronous because it uses the
same candidate gathering as normal Connection initiation. Candidate
gathering itself is asynchronous because it may
involve DNS resolution.

CTaps notifies applications of errors
through either a return code or an error callback.
Return codes are used for synchronous errors, while callbacks
are used for asynchronous errors.
Error codes in CTaps are
negative values matching
an errno corresponding to the given error. This matches the
error convention used by libuv~\cite{libuv_errors}.

In theory, error handling could be unified into using callbacks
for all errors, but CTaps splits it into two categories in order
to fail as early as possible and to differentiate simple errors,
such as illegal parameters, from more advanced asynchronous problems,
which may be outside an application's control.

As described in \cref{sec:dependency-libuv}, CTaps
uses a single global event loop internally.
This event loop must be initialized before any asynchronous operation can be initiated.
This initialization is done through the setup function \mono{ct_initialize()}, and teardown has to
be performed through \mono{ct_close()}.

After initializing CTaps, applications must perform any setup
required and either start a Listener or initiate a Connection.
After this has been performed the
application must start the event loop with \mono{ctaps_start_event_loop()}.
This event loop will block the invoking thread until there are 
no more open handles performing work.

CTaps is not designed with multi-threading in mind and does not make use
of the threading API~\cite{libuv_threading} available through libuv. CTaps expects applications
to only interact with CTaps through a single thread.

\subsection{Ownership Model}
Because all objects in CTaps are allocated on the heap,
it is important that applications free allocated objects at the appropriate times.
In order to make this as intuitive and consistent as possible, CTaps has a clear ownership model
where applications retain ownership of any allocated objects.
In addition to the objects allocated synchronously, applications
also gain ownership of Connection and Listener objects which are created asynchronously.

CTaps passes ownership of created objects to applications if
applications have not registered the
correct callbacks for receiving the object.
This will lead to a memory leak if a Connection is initiated without
an associated \mono{ready()} callback, since the Connection will never be freed.

When freeing objects which require teardown, such as Connections
or Listeners, applications must first
invoke the associated \mono{close()} function. The object
will then only be safe to free after the associated \mono{closed()}
callback has been invoked by CTaps.
Connection data and waiting messages can still be read from a connection after it has
closed, but the data is invalidated on free.

The only location where applications do not gain ownership over a passed object is in
the \mono{receive()} callbacks when receiving Messages. 
The received Message objects and Message Contexts only remain valid for the duration
of the callback, and must be deep copied if they are to be referenced outside the
callback scope. CTaps deviates from the standard model here because of the frequency
of receiving Messages and because it is assumed that Messages are handled immediately
and are rarely stored for long term use.

\section{QUIC Under TAPS} \label{sec:quicUnderTaps}
As the first open source TAPS implementation to feature QUIC as
an underlying protocol option, mapping which features of
QUIC are available through the default TAPS interface is one of the
key research questions for this thesis. Since TAPS is an interface to
a generic transport layer, TAPS concepts, terminology and intended functionality
naturally do not map one-to-one onto concepts, terminology and functionality
of QUIC. This section will describe how QUIC can be mapped to TAPS and the
related issues. A survey of which QUIC features are available
through TAPS can be seen in \cref{tab:quic-taps-mapping}. The surveyed features
in this table encompass all 17 core QUIC transport parameters defined in
RFC~9000~\cite[sec.~18.2]{rfc9000}, supplemented by higher-level protocol features such as streams,
connection migration, and the QUIC datagram and multipath extensions.

\subsection{QUIC Connections} \label{sec:quic-arch-mapping}
While the TAPS RFCs were designed with QUIC in mind,
RFC~9623 does not feature specific transport-layer protocol considerations for QUIC,
although a draft does exist~\cite{transportServicesMappingForQuic}.

The most important design decision for mapping QUIC 
to TAPS is to map a QUIC Connection
to a Connection Group instead of to a Connection.
This is natural since cloning is utilized for multistreaming, \mono{clone()}
creates new Connections and each new stream will therefore correspond to a Connection. This matches
the \enquote{stream mapping} as described for SCTP in RFC~9623~\cite[sec.~10.6]{rfc9623}
and is the recommended approach in the TAPS Mapping for QUIC
draft~\cite[sec.~1]{transportServicesMappingForQuic}.


\subsection{QUIC Streams}\label{sec:quic-streams-under-taps}
Since QUIC connections map to a TAPS Connection Group, TAPS Connections
naturally map to QUIC streams. This mapping works well and produces
behavior consistent with what is expected from the TAPS API, except
for when closing streams.

A QUIC stream is intended to transfer
a single piece of data, with both the client and the server independently marking the transfer
as completed to close the stream.
This means that the only way for a QUIC stream to close gracefully is
for a FIN flag to be sent in both directions.

As a contrast, TCP connections only consist of a single stream of data and
while they can be closed one-way, the typical behavior is that
when it is closed by one peer, the other peer closes the connection
in conjunction.

This means that a choice has to be made for how to close individual
TAPS Connections that use QUIC as the underlying protocol. Based on TCP, the 
expected behavior
is that closing an individual Connection will close both the sending
and receiving side and start teardown. For QUIC, this can in theory be emulated by sending a \mono{FIN}
flag and a \mono{STOP_SENDING} frame. The peer then must respond with a \mono{RESET_STREAM}
frame if their stream is in a ready or sending state~\cite[sec.~3.5]{rfc9000}.
A \mono{RESET_STREAM} frame tears down a QUIC stream in the sending direction,
and in combination with the sent \mono{FIN} flag, would completely close the stream.

Even though this handling of underlying QUIC streams in TAPS is
closer to the expected behavior from TCP, \mono{STOP_SENDING} is not a graceful close and
typically indicates that an application is no longer reading received data \cite[sec.~3.5]{rfc9000}.
For a graceful teardown of a single
stream only, the expected behavior is both sides sending \mono{FIN} flags when
they are done sending.

In order to avoid unnecessary error behavior for QUIC, the decision made
in CTaps is to send a \mono{FIN} flag when an application invokes \mono{close()},
letting the Connection remain open in the receiving direction.
This means an application will still have to monitor the Connection
for incoming messages, and it will not receive  a closed()-callback until
after the peer has finished sending its data.

If a peer wants to gracefully close
a connection before all data has been received, the only way of handling it gracefully is to
close the entire QUIC connection. This naturally maps to closing the entire Connection Group. 

If an application takes care and consistently interacts
with Connection Groups instead of individual Connections it can achieve the same
behavior irrespective of the underlying protocol in CTaps.
This is an example of internal functionality leaking through to
how an application should interact with the TAPS API.

The mapping of QUIC connections to TAPS Connection Groups also affects how
CTaps handles the notification of path changes. In TAPS, the \mono{path_change()} 
event is per-Connection~\cite[sec. 8.3.2]{rfc9622},
while QUIC migration is per-Connection Group.
For CTaps, this means that for each path change event every CTaps Connection is notified
separately, even though the event is always common to an
entire Connection Group. This could be mitigated by allowing registration
of callbacks on a per-Connection Group basis, but this is not featured
in the current TAPS API.

There are also features related to streams in QUIC which are not available
through a TAPS Connection.
One example of this is the one-way abort
available through \mono{STOP_SENDING}, which requests that the peer abort their
sending, and \mono{RESET_STREAM}
which aborts the local sending side~\cite[sec. 19.4--19.5]{rfc9000}.
For aborting Connections, the TAPS API presents only
\mono{abort()}~\cite[sec. 10]{rfc9622}.
TAPS Implementations cannot know whether the application
wants to abort sending, receiving, or both. The natural conclusion is
therefore to abort both, losing granular control.

Another example of QUIC stream features that are not available through
TAPS is initiating unidirectional streams. TAPS
supports setting the directionality of a Connection through the
\mono{direction} Selection Property \cite[sec.~6.2.16]{rfc9622},
but Selection Properties are only set before initiating a Connection.
This means that TAPS does not currently support
the directionality of streams when cloning.

A TAPS System is able to receive a unidirectional stream initiated by the peer as expected by QUIC.
It can infer from the incoming stream ID that a stream is unidirectional~\cite[sec. 2.1]{rfc9000}
and mark the resulting Connection as closed for sending before presenting it to the application.

\subsection{Idle Timeout}
TAPS does not feature a timeout value
which can be passed as QUIC's \mono{max_idle_timeout}~\cite[sec. 10.1]{rfc9000}
transport parameter.
TAPS does support timeouts when sending data through
the \mono{connTimeout} Connection Property~\cite[sec. 8.1.3]{rfc9622}, but
it does not support a generic idle timeout for when no data
is being transferred.

\subsection{Flow Control}
TAPS exposes flow control through enqueueing receive callbacks instead of the
maximum data parameters used by QUIC. A TAPS System can set a conservative
default for these parameters and increase them gradually as the application invokes \mono{receive()}.

\subsection{Using Multiple Paths}
QUIC path migration is supported in TAPS, and the migration
policy can be controlled through the
\mono{multipath} Selection Property.

The combination of the \mono{multipath} Selection Property and the \mono{multipathPolicy} 
Connection Property~\cite[sec. 8.1.7]{rfc9622} allows applications to express multipath preferences,
such as requesting server-initiated migration,
that go beyond what QUIC currently supports.
This is expected from TAPS,
and it is clearly documented in RFC~9622 that Connection Properties
will only have an effect if they can actually be mapped onto the
underlying connection~\cite[sec. 5]{rfc9622}. 

An overview of the expected behavior of the TAPS System based
on \mono{multipath} and \mono{multipathPolicy} can be seen in
\cref{tab:multipath-combinations}.
The \mono{interactive} and \mono{aggregate} \mono{multipath} policies require
simultaneous use of multiple paths, which is not currently supported in QUIC.
The multipath extension for QUIC~\cite{ietf-quic-multipath-21} is currently
being standardized, and it would enable these policies to be fully utilized.

\begin{table}[htbp]
  \centering
  \begin{threeparttable}
      \caption[Behavior of \mono{multipath}/\mono{multipathPolicy} in TAPS]{Effective behavior for \mono{multipath}/\mono{multipathPolicy} combinations.}
    \label{tab:multipath-combinations}
    \begin{tabular}{llp{7cm}}
      \toprule
      \rowcolor{gray!25}
      \textbf{\mono{multipath}} & \textbf{\mono{multipathPolicy}} & \textbf{Description} \\
      \midrule
      \multicolumn{3}{l}{\textbf{\texttt{disabled}}\tnote{1}} \\
        & \multicolumn{1}{|l||}{---}
        & \mono{multipathPolicy} has no effect if multipath is disabled. \\
      \midrule
      \multicolumn{3}{l}{\textbf{\texttt{active}}} \\
        & \multicolumn{1}{|l||}{\texttt{handover}}    & Will initiate migration on path failure. \\[8pt]
        & \multicolumn{1}{|l||}{\texttt{interactive}} & Will use multiple paths concurrently to minimize latency. \\[8pt]
        & \multicolumn{1}{|l||}{\texttt{aggregate}}   & Will use multiple paths concurrently to maximize throughput. \\[8pt]
      \midrule
      \multicolumn{3}{l}{\textbf{\texttt{passive}}} \\
        & \multicolumn{1}{|l||}{\texttt{handover}}    & Will accept migration and multipath used by the remote. \\[8pt]
        & \multicolumn{1}{|l||}{\texttt{interactive}} & Will accept migration and multipath from the remote, will attempt to minimize latency.\\[8pt]
        & \multicolumn{1}{|l||}{\texttt{aggregate}}   & Will accept migration and multipath from the remote, will attempt to maximize throughput.\\[8pt]
      \bottomrule
    \end{tabular}
    \begin{tablenotes}
      \footnotesize
  \item[1] For QUIC servers, which cannot migrate~\cite[sec. 1]{rfc9000}, \mono{disabled}
           can be used to set the \mono{disable_active_migration} transport parameter.
    \end{tablenotes}
  \end{threeparttable}
\end{table}

TAPS does not directly expose an equivalent to the
\mono{preferred_address} 
transport parameter~\cite[sec. 18.2]{rfc9000}, but a
value can be inferred based on the \mono{interface}
and \mono{advertisesAltaddr} Selection Properties.
If \mono{advertisesAltaddr} was set to true then
a TAPS System could infer that the application would like
incoming connections to migrate to its highest
preference interface, and
set \mono{preferred_address} accordingly.

\begin{table}[htbp]
\centering
\begin{threeparttable}
    \caption[QUIC feature mapping to TAPS API]{Mapping of QUIC features to the TAPS API.}
\label{tab:quic-taps-mapping}
\small
\begin{tabularx}{\textwidth}{lX}
\toprule
\rowcolor{gray!25}
\textbf{QUIC Feature} & \textbf{TAPS Mapping} \\
\midrule
\multicolumn{2}{l}{\textit{Available through TAPS API}} \\
\midrule
\multicolumn{1}{l||}{0-RTT data} & \mono{init_with_send()}\\
\multicolumn{1}{l||}{ALPN negotiation} & Security Parameter \\
\multicolumn{1}{l||}{Bidirectional streams} & \mono{clone()}\\
\multicolumn{1}{l||}{Connection Migration\tnote{1}} & Multiple Local Endpoints \& Selection Property\\
\multicolumn{1}{l||}{\mono{disable_active_migration}} & Selection Property\\
\multicolumn{1}{l||}{ECN} & Message Property \\
\multicolumn{1}{l||}{\mono{initial_max_streams_bidi}} & Connection Property \\
\multicolumn{1}{l||}{Key Update~\cite[sec. 6]{rfc9001}} & Security Parameter Callback \\
\multicolumn{1}{l||}{\mono{max_udp_payload_size}\tnote{2}} & Connection Property \\
\multicolumn{1}{l||}{Multipath QUIC} & Multiple Endpoints \& Selection Property\\
\multicolumn{1}{l||}{PING Frames\tnote{3}} & Connection Property \\
\multicolumn{1}{l||}{\mono{preferred_address}\tnote{4}} & Selection Properties\\
\multicolumn{1}{l||}{QUIC connection\tnote{5}} & Connection Group\\
\multicolumn{1}{l||}{QUIC datagrams~\cite{rfc9221}} & Message Property \\
\multicolumn{1}{l||}{QUIC stream} & Connection\\
\multicolumn{1}{l||}{Stream Priority} & Connection Property\\
\multicolumn{1}{l||}{TLS certificates} & Security Parameter\\
\multicolumn{1}{l||}{Unidirectional stream close (FIN)} & \mono{close()} / Message Property\\
\multicolumn{1}{l||}{Flow Control} & Scheduling \mono{receive()} invocations\\
\midrule
\multicolumn{2}{l}{\textit{Partially available through TAPS API}} \\
\midrule
\multicolumn{1}{l||}{\mono{STOP_SENDING}\tnote{6}} & \mono{abort()}\\
\multicolumn{1}{l||}{\mono{RESET_STREAM}\tnote{6}} & \mono{abort()}\\
\multicolumn{1}{l||}{\mono{ack_delay_exponent}\tnote{7}} & Connection Property\\ 
\multicolumn{1}{l||}{\mono{max_ack_delay}\tnote{7}} & Connection Property \\ 
\multicolumn{1}{l||}{Unidirectional streams\tnote{8}} & Connection Property \\
\midrule
\multicolumn{2}{l}{\textit{Not available through TAPS API}} \\
\midrule
\multicolumn{1}{l||}{Idle Timeout\tnote{9}} & ---\\
\multicolumn{1}{l||}{Directional \mono{initial_max_streams_*}} & ---\\
\multicolumn{1}{l||}{Stateless reset} & --- \\
\multicolumn{1}{l||}{Connection ID handling\tnote{10}} & ---\\
\bottomrule
\end{tabularx}
\begin{tablenotes}
\footnotesize
\item[1] Abstracted from the application, but supported through multiple Local Endpoints and \mono{multipath}.
\item[2] Requires that \mono{recvMsgMaxLen} \cite[sec. 8.1.11.5]{rfc9622} is communicated to the peer.
\item[3] Through the \mono{keepAliveTimeout} Connection Property.
\item[4] Inferred from a combination of \mono{advertisesAltAddr} and \mono{interface} Selection Properties.
\item[5] A QUIC connection maps to a Connection Group rather than a single Connection, since TAPS has no concept of streams within a Connection.
\item[6] Only available in unison through TAPS, not individually.
\item[7] Coarse control is available by setting the \mono{connCapacityProfile} Connection Property.
\item[8] TAPS can receive unidirectional streams and set the \mono{canSend} Connection Property automatically,
    but there is no support for marking the underlying stream as unidirectional when cloning.
\item[9] TAPS can set a timeout, but only for \enquote{trying to reliably deliver data to the Remote Endpoint}~\cite[sec. 8.1.3]{rfc9622}.
\end{tablenotes}
\end{threeparttable}
\end{table}

\section{Candidate Racing}
Candidate racing maps well to connection-oriented
protocols, where an endpoint can be confirmed to be reachable
by the handshake completing. It does however not map
well onto connectionless protocols such as UDP, where
it is not necessarily possible to verify that an endpoint is reachable.
A design decision therefore had to be made in CTaps on how to perform
racing when UDP was among the viable protocols.

Since UDP has no handshake to complete, CTaps has chosen to assume that a Connection
using UDP is established as soon as a socket has been opened. This means that
if the Remote Endpoint does not support UDP, then candidate racing will be unable to detect
this and UDP can win the race, without being supported by the remote.

A possible mitigation for this would be through \term{transport options}
as described in RFC~9868~\cite{rfc9868}.
Transport options are a backwards compatible way of extending UDP functionality.
Two of these options are the echo request
and the echo response~\cite[sec.~11.7]{rfc9868}.
The echo request sends a 4-byte token requesting that the peer
respond with the same token.

For candidate racing in CTaps, echo requests and echo
responses could be used as a replacement for a UDP handshake. There are however two
large issues with utilizing transport options like this.
First of all, there is no way to differentiate between the peer not being reachable
by UDP and the peer simply not supporting this transport option.
Secondly, the use of transport options is the responsibility of the layer
utilizing UDP, not the UDP layer itself~\cite[sec.~11.7]{rfc9868}. This means that if a server
is reachable over UDP, then CTaps would be reliant on the server application
itself supporting this transport option, not only the network stack.

Since the default TAPS Selection Properties rank UDP lower than other supported protocols,
the risk of UDP being ranked higher than other protocols when
the endpoint does not support UDP was deemed to be quite low. The simple solution
of treating a Connection as established as soon as the underlying UDP socket has
been opened was therefore deemed acceptable.

\section{Security}\label{sec:design-sec}
While the security of a TAPS System is primarily reliant on the security of the underlying protocols,
any implementation of TAPS needs to take security into account in its design.

The largest security risk for an application using a TAPS System is that
the required security would be silently downgraded. This risk is increased
because insecure protocols have an inherent advantage in candidate racing since they
require fewer RTTs to establish.
TAPS Systems therefore need to mandate security by default and not race
secure and insecure protocols.

This is addressed in RFC~9622, which specifies that Connections
using TAPS should use security in general, but that applications
can initialize disabled or opportunistic Security Parameters
to remain compatible with endpoints that do not support transport
security protocols~\cite[sec.~6.3]{rfc9622}.

With disabled Security Parameters, a TAPS System will not attempt
to establish any transport security, while opportunistic Security Parameters
allow for the coexistence of secure and insecure transport-layer protocols.
When racing secure and insecure protocols, TAPS implementations
should take the increased number of RTTs required to establish a secure protocol
into account. This could take the form of a longer stagger delays for insecure
protocols.

Since TAPS is a generic API,
there may be situations where the specific implementation of a TAPS System
places a limitation on the security of an application.
For example, CTaps does not currently support keeping key
files in memory and can only refer to them through filenames. This could be
seen as a CTaps-specific security gap, depending on an application's requirements.
The TAPS API itself only specifies that server and client certificates should be settable through the use of
certificate bundles, but does not mandate the form which these bundles should take.

While TAPS aims to provide reasonable security defaults,
its default configuration trades some security guarantees 
for usability and performance. One example is session isolation.
The \mono{isolateSession} Security Parameter is false by default,
which allows a TAPS implementation
to use internally cached data.
Using this cache could lead to increased connection linkability. 
An observer could see
different connections exhibiting the same behavior
and infer that they belong to the same client. This is an acceptable risk
for most applications, but security-sensitive applications
must explicitly configure parameters rather than rely on TAPS
defaults.

Fingerprinting may also be an issue inherent in TAPS itself.
Using a networking library which standardizes the interaction
with the network could lead to fingerprinting
of clients. This is an inherent property of software libraries whose behavior is exposed to the internet.
Without a large and diverse user base, fairly accurate assumptions can
be made about the users behind a connection, if the underlying library being
used can be ascertained. Fingerprinting becomes a larger concern if the behavior
of the underlying library is fairly unique, such as with the extensive
candidate racing in TAPS.
